% ═══════════════════════════════════════════════════════════════
% MUSTERLÖSUNG: Waagerechter Wurf (Physik Klasse 10E)
% E-Kurs Gymnasium-Niveau (★★★)
% ═══════════════════════════════════════════════════════════════

\documentclass[10pt, a4paper]{article}

% ═══════════════════════════════════════════════════════════════
% PAKETE
% ═══════════════════════════════════════════════════════════════

\usepackage[utf8]{inputenc}
\usepackage[T1]{fontenc}
\usepackage[ngerman]{babel}
\usepackage{helvet}
\renewcommand{\familydefault}{\sfdefault}

\usepackage[margin=1.5cm, top=1.8cm, bottom=1.5cm]{geometry}
\usepackage{enumitem}
\usepackage{tabularx}
\usepackage{booktabs}
\usepackage{array}
\usepackage{multicol}
\usepackage{amsmath}
\usepackage{amssymb}

\usepackage{xcolor}
\usepackage{tcolorbox}
\usepackage{fancyhdr}
\usepackage{lastpage}
\usepackage{tikz}

% Kompaktere Listen
\setlist{nosep, leftmargin=1.5em}

% ═══════════════════════════════════════════════════════════════
% FARBEN
% ═══════════════════════════════════════════════════════════════

\definecolor{physik}{HTML}{2c5aa0}
\definecolor{grau}{HTML}{888888}
\definecolor{hellgrau}{HTML}{f0f0f0}
\definecolor{success}{HTML}{2a7a4b}
\definecolor{warning}{HTML}{b35c00}
\definecolor{loesung}{HTML}{c62828}

\newcommand{\fachfarbe}{physik}

% ═══════════════════════════════════════════════════════════════
% VARIABLEN
% ═══════════════════════════════════════════════════════════════

\newcommand{\thema}{Waagerechter Wurf}
\newcommand{\fach}{Physik}
\newcommand{\klasse}{Klasse 10E}
\newcommand{\stunde}{Bewegungslehre}

% ═══════════════════════════════════════════════════════════════
% KOPF- UND FUSSZEILE
% ═══════════════════════════════════════════════════════════════

\pagestyle{fancy}
\fancyhf{}
\renewcommand{\headrulewidth}{0.4pt}
\renewcommand{\footrulewidth}{0pt}

\lhead{\small\fach{} \klasse{} -- \thema{} \textcolor{loesung}{\textbf{(MUSTERLÖSUNG)}}}
\rhead{\small}

\cfoot{\small\thepage/\pageref{LastPage}}

% ═══════════════════════════════════════════════════════════════
% BEFEHLE
% ═══════════════════════════════════════════════════════════════

% Lösung (rot)
\newcommand{\lsg}[1]{\textcolor{loesung}{\textbf{#1}}}

% Physik-Highlight
\newcommand{\ph}[1]{\textcolor{\fachfarbe}{\textbf{#1}}}

% Sternchen für Schwierigkeit
\newcommand{\stmark}{$\star$}

% Merksatz-Box
\newtcolorbox{merksatz}[1][]{
    colback=hellgrau,
    colframe=\fachfarbe,
    boxrule=0.8pt,
    arc=0pt,
    left=6pt, right=6pt, top=6pt, bottom=6pt,
    fonttitle=\bfseries\small,
    title=#1
}

% Formel-Box
\newtcolorbox{formelbox}{
    colback=white,
    colframe=\fachfarbe,
    boxrule=1.5pt,
    arc=0pt,
    left=8pt, right=8pt, top=8pt, bottom=8pt
}

% Hinweis-Box
\newtcolorbox{hinweisbox}{
    colback=white,
    colframe=warning,
    boxrule=1pt,
    arc=0pt,
    left=6pt, right=6pt, top=6pt, bottom=6pt
}

% Abschnittstitel
\newcommand{\abschnitt}[1]{%
    \vspace{8pt}
    \noindent\textbf{\textcolor{\fachfarbe}{#1}}
    \vspace{4pt}
}

% Punkte am Rand
\newcommand{\punkte}[1]{\hfill\small\textcolor{grau}{/#1 P.}}

% ═══════════════════════════════════════════════════════════════
% DOKUMENT
% ═══════════════════════════════════════════════════════════════

\begin{document}

% ═══════════════════════════════════════════════════════════════
% TITEL
% ═══════════════════════════════════════════════════════════════

\begin{center}
{\large\bfseries Arbeitsblatt: \thema{} -- \textcolor{loesung}{MUSTERLÖSUNG}}\\[2pt]
{\small\textcolor{grau}{Lösungen in Rot}}
\end{center}

\vspace{-2pt}
\hrule
\vspace{8pt}

% ═══════════════════════════════════════════════════════════════
% ZWEISPALTIG
% ═══════════════════════════════════════════════════════════════

\begin{multicols}{2}
\small

% ─────────────────────────────────────────────────────────────
% KASTEN 1: Die Kernerkenntnis
% ─────────────────────────────────────────────────────────────

\abschnitt{Kasten 1: Das Überlagerungsprinzip}

\begin{merksatz}[Merksatz]
Beim waagerechten Wurf überlagern sich zwei \ph{unabhängige} Bewegungen:

\vspace{6pt}
\begin{tabular}{ll}
\textbf{Horizontal:} & \lsg{gleichförmige} Bewegung \\[4pt]
\textbf{Vertikal:} & \lsg{freier Fall} \\
\end{tabular}

\vspace{8pt}
Die horizontale Geschwindigkeit beeinflusst die

\lsg{Fallzeit} \textbf{nicht}.

\vspace{6pt}
Deshalb landen ein fallender und ein geworfener Ball,

die gleichzeitig starten, \lsg{gleichzeitig}.
\end{merksatz}

\vspace{4pt}

% ─────────────────────────────────────────────────────────────
% SKIZZE
% ─────────────────────────────────────────────────────────────

\abschnitt{Skizze: Flugbahn mit Geschwindigkeiten}

\begin{center}
\begin{tikzpicture}[scale=0.6]
    % Koordinatensystem
    \draw[->] (0,0) -- (7,0) node[right] {$x$};
    \draw[->] (0,0) -- (0,5) node[above] {$y$};

    % Boden
    \draw[thick] (-0.5,0) -- (7,0);

    % Plattform
    \fill[gray!50] (-0.5,4.2) rectangle (0.5,4.4);

    % Höhe h
    \draw[<->] (-0.3,0) -- (-0.3,4.2) node[midway, left] {$h$};

    % Parabel (rot für Lösung)
    \draw[thick, loesung] (0.5,4.2) parabola (6,0);

    % Ball Start
    \fill[physik] (0.5,4.2) circle (0.12);

    % Ball Mitte
    \fill[physik] (3.2,2.5) circle (0.12);

    % Ball Ende
    \fill[physik] (6,0) circle (0.12);

    % Wurfweite
    \draw[<->] (0.5,-0.5) -- (6,-0.5) node[midway, below] {$s_x$};

    % v0 Pfeil (Start)
    \draw[->, thick, loesung] (0.5,4.2) -- (1.8,4.2) node[above] {$\vec{v}_0$};

    % v_x Pfeil (Mitte) - horizontal
    \draw[->, thick, loesung] (3.2,2.5) -- (4.5,2.5) node[above] {$\vec{v}_x$};

    % v_y Pfeil (Mitte) - vertikal nach unten
    \draw[->, thick, blue] (3.2,2.5) -- (3.2,1.2) node[right] {$\vec{v}_y$};

    % v_res Pfeil (Mitte) - diagonal
    \draw[->, very thick, success] (3.2,2.5) -- (4.5,1.2) node[right] {$\vec{v}_{res}$};

    % Auftreffwinkel am Ende
    \draw[dashed, gray] (6,0) -- (6,1.5);
    \draw[thick, loesung] (6,0.8) arc (90:143:0.8);
    \node at (5.4,0.9) {\footnotesize\lsg{$\alpha$}};
\end{tikzpicture}
\end{center}

% ─────────────────────────────────────────────────────────────
% KASTEN 2: Grundformeln
% ─────────────────────────────────────────────────────────────

\abschnitt{Kasten 2: Grundformeln}

\begin{formelbox}
\textbf{Schritt 1: Fallzeit}
\vspace{2pt}

\begin{equation*}
t = \sqrt{\frac{2h}{g}}
\end{equation*}

\vspace{2pt}
\footnotesize
$t$ = Fallzeit [s] \quad $h$ = Höhe [m] \quad $g \approx 10\,\text{m/s}^2$
\end{formelbox}

\vspace{4pt}

\begin{formelbox}
\textbf{Schritt 2: Wurfweite}
\vspace{2pt}

\begin{equation*}
s_x = v_0 \cdot t
\end{equation*}

\vspace{2pt}
\footnotesize
$s_x$ = Wurfweite [m] \quad $v_0$ = Anfangsgeschw. [m/s]
\end{formelbox}

\vspace{4pt}

% ─────────────────────────────────────────────────────────────
% KASTEN 3: Erweiterte Formeln
% ─────────────────────────────────────────────────────────────

\abschnitt{Kasten 3: Erweiterte Formeln (E-Kurs)}

\begin{formelbox}
\textbf{Vertikalgeschwindigkeit beim Aufprall}
\vspace{2pt}

\begin{equation*}
v_y = g \cdot t
\end{equation*}

\vspace{2pt}
\footnotesize
$v_y$ = Vertikalgeschw. beim Aufprall [m/s]
\end{formelbox}

\vspace{4pt}

\begin{formelbox}
\textbf{Resultierende Geschwindigkeit}
\vspace{2pt}

\begin{equation*}
v_{res} = \sqrt{v_x^2 + v_y^2}
\end{equation*}

\vspace{2pt}
\footnotesize
$v_{res}$ = Aufprallgeschwindigkeit [m/s] \quad (Pythagoras!)
\end{formelbox}

\vspace{4pt}

\begin{formelbox}
\textbf{Auftreffwinkel}
\vspace{2pt}

\begin{equation*}
\tan(\alpha) = \frac{v_y}{v_x} \quad \Rightarrow \quad \alpha = \arctan\left(\frac{v_y}{v_x}\right)
\end{equation*}

\vspace{2pt}
\footnotesize
$\alpha$ = Winkel zur Horizontalen beim Aufprall
\end{formelbox}

\vspace{4pt}

\begin{hinweisbox}
\footnotesize
\textbf{Hinweis:} Wir rechnen mit $g \approx 10\,\text{m/s}^2$ für einfache Zahlen.\\
Der exakte Wert ist $g = 9{,}81\,\text{m/s}^2$.
\end{hinweisbox}

% ─────────────────────────────────────────────────────────────
% Spaltenumbruch
% ─────────────────────────────────────────────────────────────

\columnbreak

% ─────────────────────────────────────────────────────────────
% TABELLE
% ─────────────────────────────────────────────────────────────

\abschnitt{Tabelle: Typische Fallzeiten}

\begin{center}
\begin{tabular}{|c|c|}
\hline
\textbf{Höhe $h$} & \textbf{Fallzeit $t$} \\
\hline
5 m & 1 s \\
\hline
20 m & 2 s \\
\hline
45 m & 3 s \\
\hline
80 m & 4 s \\
\hline
\end{tabular}
\end{center}

\vspace{2pt}

% ─────────────────────────────────────────────────────────────
% AUFGABE 1 MIT LÖSUNG
% ─────────────────────────────────────────────────────────────

\abschnitt{Aufgabe 1: Grundaufgabe (\stmark\stmark)}\punkte{4}

Ein Ball wird aus $h = 20\,\text{m}$ Höhe waagerecht mit $v_0 = 15\,\text{m/s}$ geworfen.

\textbf{a)} Berechne die Fallzeit.

\textcolor{loesung}{$t = \sqrt{\dfrac{2h}{g}} = \sqrt{\dfrac{2 \cdot 20\,\text{m}}{10\,\text{m/s}^2}} = \sqrt{4\,\text{s}^2}$ \quad $\boxed{t = 2\,\text{s}}$}

\vspace{4pt}

\textbf{b)} Berechne die Wurfweite.

\textcolor{loesung}{$s_x = v_0 \cdot t = 15\,\text{m/s} \cdot 2\,\text{s}$ \quad $\boxed{s_x = 30\,\text{m}}$}

\vspace{4pt}

% ─────────────────────────────────────────────────────────────
% AUFGABE 2 MIT LÖSUNG
% ─────────────────────────────────────────────────────────────

\abschnitt{Aufgabe 2: Anwendung (\stmark\stmark)}\punkte{5}

Ein Flugzeug fliegt in $h = 45\,\text{m}$ Höhe mit $v_0 = 20\,\text{m/s}$. Es lässt ein Hilfspaket fallen.

\textbf{a)} Wie lange ist das Paket in der Luft?

\textcolor{loesung}{$t = \sqrt{\dfrac{2 \cdot 45\,\text{m}}{10\,\text{m/s}^2}} = \sqrt{9\,\text{s}^2}$ \quad $\boxed{t = 3\,\text{s}}$}

\vspace{4pt}

\textbf{b)} Wie weit vor dem Ziel muss abgeworfen werden?

\textcolor{loesung}{$s_x = 20\,\text{m/s} \cdot 3\,\text{s}$ \quad $\boxed{s_x = 60\,\text{m}}$}

\vspace{4pt}

% ─────────────────────────────────────────────────────────────
% AUFGABE 3 MIT LÖSUNG
% ─────────────────────────────────────────────────────────────

\abschnitt{Aufgabe 3: Erweitert (\stmark\stmark\stmark)}\punkte{8}

Ein Klippenspringer springt aus $h = 80\,\text{m}$ Höhe mit $v_0 = 30\,\text{m/s}$ ab.

\textbf{a)} Berechne die Fallzeit $t$.

\textcolor{loesung}{$t = \sqrt{\dfrac{2 \cdot 80}{10}}\,\text{s} = \sqrt{16}\,\text{s} = \boxed{4\,\text{s}}$}

\vspace{2pt}

\textbf{b)} Berechne die Wurfweite $s_x$.

\textcolor{loesung}{$s_x = 30\,\text{m/s} \cdot 4\,\text{s} = \boxed{120\,\text{m}}$}

\vspace{2pt}

\textbf{c)} Berechne die Vertikalgeschwindigkeit $v_y$ beim Aufprall.

\textcolor{loesung}{$v_y = g \cdot t = 10\,\text{m/s}^2 \cdot 4\,\text{s} = \boxed{40\,\text{m/s}}$}

\vspace{2pt}

\textbf{d)} Berechne die Aufprallgeschwindigkeit $v_{res}$.

\textcolor{loesung}{$v_{res} = \sqrt{v_x^2 + v_y^2} = \sqrt{30^2 + 40^2}\,\text{m/s} = \sqrt{2500}\,\text{m/s}$ \quad $\boxed{v_{res} = 50\,\text{m/s}}$ (3-4-5-Dreieck!)}

\vspace{4pt}

% ─────────────────────────────────────────────────────────────
% AUFGABE 4 MIT LÖSUNG
% ─────────────────────────────────────────────────────────────

\abschnitt{Aufgabe 4: Transfer (\stmark\stmark\stmark)}\punkte{6}

Ein Stuntman springt von einem Dach und landet $s_x = 90\,\text{m}$ entfernt mit einer Aufprallgeschwindigkeit von $v_{res} = 50\,\text{m/s}$.

\textbf{a)} Berechne $v_0$, wenn $v_y = 40\,\text{m/s}$.

\textcolor{loesung}{Pythagoras rückwärts: $v_{res}^2 = v_0^2 + v_y^2$ $\Rightarrow$ $v_0 = \sqrt{v_{res}^2 - v_y^2} = \sqrt{50^2 - 40^2}\,\text{m/s} = \sqrt{900}\,\text{m/s}$ \quad $\boxed{v_0 = 30\,\text{m/s}}$}

\vspace{2pt}

\textbf{b)} Berechne die Fallzeit $t$.

\textcolor{loesung}{Umstellen: $s_x = v_0 \cdot t \Rightarrow t = \dfrac{s_x}{v_0}$ \quad $t = \dfrac{90\,\text{m}}{30\,\text{m/s}} = \boxed{3\,\text{s}}$}

\end{multicols}

% ═══════════════════════════════════════════════════════════════
% SELBSTCHECK
% ═══════════════════════════════════════════════════════════════

\vspace{4pt}
\hrule
\vspace{6pt}

\small
\textbf{Selbstcheck -- Ich kann...}

\begin{tabular}{|p{6cm}|c|c||p{6cm}|c|c|}
\hline
\textbf{Kompetenz} & \textbf{Ja} & \textbf{Nein} & \textbf{Kompetenz} & \textbf{Ja} & \textbf{Nein} \\
\hline
...erklären, warum beide Bälle gleichzeitig landen. & \lsg{$\boxtimes$} & $\square$ &
...die Aufprallgeschwindigkeit $v_{res}$ berechnen. & \lsg{$\boxtimes$} & $\square$ \\
\hline
...die Fallzeit aus der Höhe berechnen. & \lsg{$\boxtimes$} & $\square$ &
...den Auftreffwinkel $\alpha$ berechnen. & \lsg{$\boxtimes$} & $\square$ \\
\hline
...die Wurfweite berechnen. & \lsg{$\boxtimes$} & $\square$ &
...reale Anwendungen analysieren. & \lsg{$\boxtimes$} & $\square$ \\
\hline
\end{tabular}

% ═══════════════════════════════════════════════════════════════
% PUNKTE-ÜBERSICHT
% ═══════════════════════════════════════════════════════════════

\vfill
\hrule
\vspace{4pt}
\small
\begin{tabular}{llllll}
\textbf{Aufgabe 1:} & \lsg{4} / 4 P. &
\textbf{Aufgabe 2:} & \lsg{5} / 5 P. &
\textbf{Aufgabe 3:} & \lsg{8} / 8 P. \\[4pt]
\textbf{Aufgabe 4:} & \lsg{6} / 6 P. &
\textbf{Gesamt:} & \lsg{23} / 23 P. &
\textbf{Note/NP:} & \lsg{15 NP / Note 1}
\end{tabular}

\end{document}
